\documentstyle[%


righttag%


,11pt%


,amssymb%


,russian%               remove in Eng.


,izvrutre%                 Set izven% in Eng.


]{amsart}





% 


% sidorov@math.isu.runnet.ru


\newcommand{\la}{\langle}


\newcommand{\ra}{\rangle}


\newcommand{\supp}{\operatorname{supp}}


\newcommand{\tts}[1]{}





%\hoffset=-15mm%                  For Eng.


\setcounter{page}{59}


\god{2001}


\nomer{9~(472)} %                        Remove in Eng.


%\nomer{Vol.\,45, No.\,7}%               Set in Eng.


%\str{\pageref{begin}--\pageref{end}}%  Set in Eng.


\udk{517.988.67}





\author{\it Н.А.\,Сидоров}


% NB:   ^^ remove in Eng.


\title{Параметризация простых разветвляющихся решений


полного ранга и итерации в нелинейном анализе}


\thanks{Работа выполнена при финансовой поддержке INTAS, грант 2000-15.}





\date{}


%\pagestyle{myheadings}%         Set in Eng.


\pagestyle{plain} %              Remove in Eng.


\begin{document}


%\thispagestyle{plain}%          Set in Eng.


\maketitle


\label{begin}








\section*{Введение}





Хорошо известна важная роль геометрии плоскостей, диаграмм и


многогранников


Ньютона [1], [2] в асимптотическом анализе задач теории ветвления


([1]--[5] и др.). При построении разветвляющихся решений итерационными


методами использовались явная и неявная параметризации [6]--[12], а в


сложных


случаях [10], [13]--[15] --- указанные выше


геометрические конструкции. В [6], [7], [9]--[12] решающее


значение для сходимости методов имели жесткие условия,


обеспечивающие выход на главную часть уравнения разветвления (у\,.р.)


уже на первой итерации. В частных случаях указывался способ построения


начального приближения и выбора параметризации искомой ветви. В


[6], [7], [9], [11], [12]  параметр униформизации выбирался


непосредственно по


виду диаграмм Ньютона коэффициентов проекции $QF$. Поэтому на каждой


итерации решалось одно  линейное уравнение. Итерационные


методы для более сложных ситуаций в [6]--[12] не были разработаны, т.\,к.


неясно было, как выбирать параметр униформизации ветвей и организовать


итерационный процесс в общем случае.





В [13], [14] показано, что существенное расширение класса задач,


создание более гибких алгоритмов возможны за счет привлечения


$N$-ступенчатых методов, когда на каждой итерации решается $N$ линейных


уравнений и применяются разные параметры униформизации согласно [2].


В [13] предложен $N$-ступенчатый итерационный метод, сходящийся в


окрестности точки ветвления при более слабых условиях, чем другие методы.


Метод обладает достоинствами явной и неявной параметризаций [12],


т.\,к. дает определенную свободу выбора параметра униформизации.


Однако итерационная схема, предложенная на его основе в ([13], \S\,2),


в общем случае требует предварительного построения элемента $\wh{x}^0$


из подпространства $E^{\infty-n}_1,$ используемого в начальном


приближении.


Способ построения элемента $\wh{x}^0$ дан в ([13], леммa 2). Тем


не менее, предварительное вычисление $\wh{x}^0$ усложняет вычисления и


саму итерационную схему.





Целью данной работы является не только модификация и расширение


возможностей итерационных методов работ [11], [12], [16], но и вывод


итерационных формул, удобных для теории приближенных методов и ее


приложений. Предложена более удобная итерационная схема с


упрощенным выбором начального приближения.








\section{Уравнение разветвления и выбор начального приближения}





В параграфе приведены необходимые


сведения из [1], [2], [12], [13] по выбору параметра


униформизации и начального приближения.





Пусть $E_1$, $E_2$ --- банаховы пространства. Рассмотрим уравнение


\begin{equation}


Bx - R(x,\lambda) = 0,


\end{equation}


где $B: D(B)\subset E_1\rightarrow E_2$ --- замкнутый


фредгольмов оператор с плотной в $E_1$ областью определения,


$\dim N(B) = n \ge 1,$ оператор


$R(x,\lambda) = R_{01}\lambda +


\sum\limits_{i+k\ge 2}R_{ik}(x)\lambda^k$ аналитический в окрестности


точки $x=0$, $\lambda =0$.


Требуется построить решение  $x\rightarrow 0$ при $\lambda


\rightarrow 0$.





Пусть $\{\varphi_i\}^n_1$ --- базис в $N(B)$, $\{\psi_i\}^n_1$ --- базис в


$N(B^*)$, $\{\gamma_i\}^n_1$, $\{z_i\}^n_1$ --- соответствующие


биортогональные системы из $E^*_1$ и $E_2$,


$P = \sum\limits_1^n\la\cdot,\gamma_i\ra\varphi_i$,


$Q = \sum\limits_1^n\la\cdot,\psi_i\ra z_i$, $\Gamma = \Big(B +


\sum\limits_1^n\la\cdot,\gamma_i\ra z_i\Big)^{-1}$.


Ограниченный оператор $\Gamma$


называют псевдорезольвентой фредгольмова оператора $B$ [1].





Полагая в уравнении (1)


\begin{equation}


x = \xi\varphi + \Gamma y,


\end{equation}


где


$$


Qy = 0, \quad \xi\varphi = \sum^n_1\xi_i\varphi_i,


$$


получим для определения $\xi$ и $y$ систему


\begin{gather}


y = R(\xi\varphi + \Gamma y, \lambda), \\


\la y,\psi_i\ra = 0, \quad i=1,\dotsc, n.


\end{gather}


На основании теоремы о неявных операторах [4] уравнение (3) имеет


единственное малое решение


\begin{equation}


y = \sum_{m\ge 2}y_{m0}(\xi\varphi) + \sum_{m\ge 0}\sum_{n\ge


1}y_{mn}(\xi\varphi)\lambda^n.


\end{equation}


Рекуррентные формулы для вычисления коэффициентов $y_{mn}$ построены в


[13].


Здесь приведем другую, более экономную в вычислениях рекуррентную формулу


$$


y_{n-j,j} = \frac{1}{(n-j)!j!}


\frac{\partial^n}{\partial\lambda^j\partial\mu^{n-j}}


\sum^n_{i+k=2}\lambda^kR_{ik}\bigg(\sum^{n+1-i-k}_{r+s=1}\Gamma


y_{rs}\mu^r\lambda^s\bigg)\bigg|_{\mu=0,\lambda=0},


$$


$j=0,1,\dotsc,n$, $n=2,3,\dotsc$, где $y_{01} = R_{01}$, $y_{01}


\triangleq \xi_1 z_1 + \dotsb + \xi_n z_n$.


Подставляя решение (5) в (4), получим уравнение разветвления


\begin{equation}


L^j(\xi,\lambda) \overset{\mathrm{def}}{=} \sum_{m\ge 2}L_{mo}^j(\xi)\lambda +


\sum_{m\ge 0}\sum_{\nu\ge 1}L_{m\nu}^j(\xi)\lambda^{\nu} = 0,\quad


j=1,\dotsc,n,


\end{equation}


где


$$


L_{m\nu}^j = \la y_{m\nu}(\xi\varphi),\psi_j\ra =


\sum_{m_1+\dotsc +m_n=m}L^j_{m_1,\dotsc,m_n,\nu}\xi_1^{m_1}


\dotsc\xi_n^{m_n}.


$$





Пусть $L(\xi_1,\dotsc,\xi_{n+1}) = \sum_i L_i\xi^i$,


$\xi^i=\xi_1^{i_1}\dotsc\xi_{n+1}^{i_{n+1}}$, ---


одна из левых частей системы (6). Для симметрии в обозначениях положено


$\lambda = \xi_{n+1}$. Не ограничивая общности, считаем, что


$L(\xi_1,\dotsc,O_i,\dotsc,\xi_{n+1}) \ne 0$ при $i=1,\dotsc,n+1$. В


противном случае мы бы сократили соответствующее уравнение в (6) на


некоторую степень $\xi_i$.





Пусть $\supp L = \{i\mid i\in N_+^{n+1},\ L_i \ne 0\}$,


$N_+$ --- множество целых положительных чисел.





\begin{defn}


Гиперплоскость


$$


l:(\xi \in R_+^{n+1}\mid (\xi,\alpha) = \Theta),\ \; \alpha\in N_+^{n+1},


\ \; \Theta\in N_+,


$$


назовем опорной для $\supp L$, если:


\begin{itemize}


\item[1)] $(\xi,\alpha) \ge \Theta $ при $\xi\in \supp L$;


\item[2)] $l\cap \supp L \ne 0$.


\end{itemize}


\end{defn}





{\bf Условие 1} [13]. {\it В $N_+$ найдены числа


$\alpha_1,\dotsc,\alpha_{n+1}$,


$\Theta_1,\dotsc,\Theta_n$ такие, что параллельные гиперплоскости


$l_j: (\xi\in R_+^{n+1}\mid (\xi,\alpha) = \Theta_j)$, $j=1,\dotsc,n$,


являются опорными соответственно для $\supp L^j$, $j=1,\dotsc,n$.}


\medskip





Аналитически условие 1 означает, что при $\xi_i =


\varepsilon^{\alpha_i}\eta_i$, $i=1,\dotsc,n+1$,


и $\varepsilon\rightarrow 0$


\begin{equation}


L^j = \varepsilon^{\Theta_j}l_j(\eta_1,\dotsc,\eta_{n+1}) +


r_j(\eta,\varepsilon),


\end{equation}


где


$$


l_j = \sum_{(i,\alpha)=\Theta_j}L_i^j\eta^i\not\equiv 0, \quad


r_j=o(\varepsilon^{\Theta_j}), \quad j=1,\dotsc,n.


$$


Вектор-функцию $l(\eta)$ называют главной частью у.\,р.





Проблема выбора вектора $\alpha$ решена в [2]. Если $\alpha_1 =\dotsb


=\alpha_m$, $\alpha_{m+1}=\dotsb =\alpha_{n+1}$, то гиперплоскости $l_j$


симметричны относительно осей $\xi_1,\dotsc,\xi_m$  и осей


$\xi_{m+1},\dotsc,\xi_{n+1}$.   В такой симметричной ситуации векторы


$\alpha$  и $\Theta$  легко строятся методом диаграмм Ньютона [1], [16].


\medskip





{\bf Условие 2} [13]. {\it Система


\begin{equation}


l_j(\eta_1,\dotsc,\eta_{n+1}) = 0, \quad j=1,\dotsc,n,


\end{equation}


имеет решение $\eta^0 = (\eta^0,\dotsc,\eta^0_{n+1})$,


причем $\det\big\|\frac{\partial


l_j(\eta^0)}{\partial\eta_i}\big\|_{j=\overline{1,n},


\ i=\overline{1,n+1}\,\setminus *} \ne 0$.}


\medskip





Решение $\eta^0 \ne 0$ назовем решением полного ранга системы (8), а


соответствующую пару $(x(\varepsilon)$, $\lambda(\varepsilon))$,


удовлетворяющую (1), простым решением [17] полного ранга уравнения (1).





Пусть выполнены условия 1, 2. Тогда на основании теоремы о неявной функции


у.\,р. (6) имеет решение


$$


\xi_i = \varepsilon^{\alpha_i}(\eta^0_i + o(1)),\ \; \xi_* =


\varepsilon^{\alpha_*}\eta^0_*,\ \; i=(1,\dotsc,n+1)\setminus *,


$$


где $\xi_{n+1}\triangleq \lambda(\varepsilon)$.


Подставляя это решение в (2) с учетом (5), получим искомое решение полного


ранга $x(\varepsilon)\rightarrow 0$, $\lambda(\varepsilon)\rightarrow 0$ при


$\varepsilon \rightarrow 0$ уравнения (1).





Вектор $\eta^0 = (\varepsilon^{\alpha_1}\eta^0_1,\dotsc,


\varepsilon^{\alpha_{n+1}}\eta^0_{n+1})$ будет использован в следующем


параграфе при построении начального приближения этого решения в


итерационном методе.





\section{Итерационная схема}





Пусть выполнены условия 1, 2. Решение уравнения (1) будем искать в виде


\begin{equation}


x = \sum_1^n\xi_i\varphi_i + \varepsilon^r\Gamma y, \quad \lambda =


\varepsilon^{\alpha_{n+1}}\eta_{n+1}(\varepsilon)


\triangleq \xi_{n+1}


\end{equation}


с условием $Qy = 0$. Здесь $\xi_i =


\varepsilon^{\alpha_i}\eta_i(\varepsilon)$,


$i=1,\dotsc,n$, $r = \min (\alpha_1,\dotsc,\alpha_{n+1})$, $\eta_i(0) =


\eta_i^0$, $\eta_*(\varepsilon) \equiv \eta^0_*$, $*\in(1,\dotsc,n+1)$,


$$


y(0) =


\begin{cases}


0, & r<\alpha_{n+1};\\


R_{01}\eta_{n+1}(0), & r = \alpha_{n+1}.


\end{cases}


$$





Отметим, что представление (9) удобнее использованного ранее (см. (19) в


[13]), т.\,к. не содержит дополнительного члена $\wh{x}^0$,


усложнявшего метод.


Неизвестные $y(\varepsilon)$ и $\eta_i(\varepsilon)$, $i =


\overline{1,n+1}\setminus *$, непрерывные в нуле, определим из системы


\begin{gather}


\varepsilon^r y = R(\xi\varphi + \varepsilon^r\Gamma y,\; \xi_{n+1})


\triangleq \Phi(y), \\


%


\la y,\psi_i\ra = 0, \quad i=1,\dotsc,n.


\end{gather}





Отметим, что


$\lim\limits_{\varepsilon\to 0}\varepsilon^{-r}\Phi(y) = y(0)$, $\eta_* =


\eta^0_*$.


Следуя [13], преобразуем систему (10), (11) так, чтобы выполнялись условия


теоремы о неявном операторе [4] в окрестности точки $y_0 = y(0)$, $\eta^0_i


= \eta_i(0)$, $i =\overline{1,n+1}\,\setminus *$.





С этой целью, введя итерации в систему (10), (11), перейдем к системе


\begin{gather}


y =\varepsilon^{-r} \underbrace{\Phi(\Phi\dotsc(\Phi(y))}_{N}\dotsc)


\overset{\mathrm{def}}{=}{} \varepsilon^{-r}\Phi^N(y),\\


%


\varepsilon^{-\Theta_j}


\bigg\la R\bigg(\sum^n_{i=1}\varepsilon^{\alpha_i}\eta_i\varphi_i


+ \Gamma\Phi^N(y), \varepsilon^{\alpha_{n+1}}\eta_{n+1}\bigg),


\Psi_j\bigg\ra  = 0,\quad


j=\overline{1,n}.


\end{gather}


Здесь $\eta_* = \eta^0_*$,


$$


\Phi^N(y) \overset{\mathrm{def}}{=}{} R(\xi\varphi + \Gamma R(\xi\varphi +


\dotsb + \Gamma R(\xi\varphi + \varepsilon^r\Gamma y,


\xi_{n+1}),\dotsc,)\xi_{n+1}).


$$


Любое решение $(y,\eta_1,\dotsc,\eta_{*-1},\eta_{*+1},\dotsc,\eta_{n+1})$


системы (12), (13) удовлетворяет системе (10), (11) и условию $Qy = 0$.





Систему (12), (13) рассмотрим как одно операторное уравнение


\begin{equation}


K(u,\varepsilon) = 0,


\end{equation}


где $K: Y\times R^1\rightarrow Y$, $Y = E_2 \oplus R^n$,


$\|u\|\triangleq \|y\| + |\eta|$,


$u = (y(\varepsilon),\eta_1(\varepsilon),\dotsc,


\eta_{*-1}(\varepsilon),\dotsc,


\eta_{n+1}(\varepsilon))'$,


$\eta_*(\varepsilon) \equiv \eta^0_*$.





На основании условий 1, 2 при достаточно большом $N \ge N_0$  оператор


$K(u,\varepsilon)$ будет непрерывным в окрестности точки $(u_0,0)$,


где $u_0 = (y(0),\eta^0)$, 


$K(u_0,0) = 0$. Используя разложение правой части формул (12), (13) в ряды


Тейлора в окрестности точки $\varepsilon = 0$, можно показать, что


достаточно взять


$N_0 = \max\limits_{1\le j\le n}(\Theta_j


- \min(\alpha_1,\dotsc,\alpha_n))$. Эта


оценка является грубой и в более конкретной ситуации легко может быть


улучшена. Смотри, например, ниже условие 3, при котором можно взять $N_0 =


1$.





При $N\ge N_0$ оператор $K(u,\varepsilon)$ имеет производную Фреше


$K_u(u,\varepsilon)$, непрерывную в точке $(u_0,0)$,


$$


K_u(u_0,0) =


\begin{pmatrix}


J & 0 \\


0 & D


\end{pmatrix}.


$$


Здесь


$$


D = \bigg\|\frac{\partial l_j(\eta_0)}{\partial\eta_i}\bigg\|


\begin{Sb}


j=\overline{1,n}\\


i=\overline{1,n+1}\,\setminus *


\end{Sb}.


$$





В силу условия 2\; $\det D\ne 0$ и, следовательно, существует ограниченный


обратный оператор


$$


K_u^{-1}(u_0,0) =


\begin{pmatrix}


J & 0 \\


0 & D^{-1}


\end{pmatrix}.


$$


Поэтому уравнение (14) удовлетворяет всем условиям теоремы о неявном


операторе [4].





Искомое решение $u = (y,\eta)$ можно найти методом последовательных


приближений


\begin{gather}


y_m = \varepsilon^{-r}R(\xi^{m-1}\varphi + \Gamma R(\xi^{m-1}\varphi +


\dotsb + \Gamma R(\xi^{m-1}\varphi + \varepsilon^r\Gamma y_{m-1},


\xi^{m-1}_{n+1}),\dotsc,) \xi^{m-1}_{n+1}),\\


%


\eta^m = \eta^{m-1} - D^{-1}\la R(\xi^{m-1}\varphi + \varepsilon^r\Gamma y_m,


\xi^{m-1}_{n+1}),E(\varepsilon)\psi\ra,\ \; m=1,2,\dotsc.


\end{gather}


Здесь $E\psi \triangleq(\varepsilon^{-


\Theta_1}\psi_1,\dotsc,\varepsilon^{-\Theta_n}\psi_n)'$,


$\xi\triangleq(\varepsilon^{\alpha_1}\eta_1,\dotsc,\varepsilon^{\alpha_n}


\eta_n)'$,


$\xi_{n+1} = \varepsilon^{\alpha_{n+1}}\eta_{n+1}$, $\eta_* \equiv


\eta^0_*$.


Начальное приближение $(y_0,\eta^0)$ определяется формулами


$$


y_0 = y(0), \quad \eta^0 = (\eta^0_1,\dotsc,\eta^0_{*-


1},\eta^0_{*+1},\dotsc,\eta^0_{n+1})',


$$


где $\eta^0$ удовлетворяет системе (8), ($'$) --- знак


транспонирования.





Преобразуем (15) к более удобному виду. Введем последовательности


\begin{gather}


x_m = \xi^m\varphi + \varepsilon^r\Gamma y_m, \\


%


\lambda_m = \xi^m_{n+1} \triangleq


\varepsilon^{\alpha_{n+1}}\eta^m_{n+1},


\end{gather}


где $\xi^m\varphi=


\sum\limits^n_{i=1}\varepsilon^{\alpha_i}\eta^m_i\varphi_i$,


$m=1,2,\dotsc$ Введем еще вспомогательные элементы


\begin{gather}


x^0_{m-1} = \xi^{m-1}\varphi + \varepsilon^r\Gamma y_{m-1}


\triangleq x_{m-1},\\


%


x^i_{m-1} = \xi^{m-1}\varphi + \Gamma R(x_{m-1}^{i-1},\lambda_{m-1}), \quad


i=1,\dotsc,N-1.


\end{gather}


Тогда итерационная схема (15), (16) примет вид


\begin{equation}


\begin{split}


y_m &= \xi^{-r}R(x_{m-1}^{N-1},\lambda_{m-1}), \\


\eta^m &= \eta^{m-1} - D^{-1}\la R(\xi^{m-1}\varphi +


\varepsilon^r\Gamma y_m,\lambda_{m-


1}),E(\varepsilon)\Psi\ra, \quad m=1,2,\dotsc


\end{split}


\end{equation}





В качестве начального приближения следует взять $y^0 = y(0)$,


$\xi_i^0 = \varepsilon ^{\alpha_i}\eta_i^0$,


$i = \overline{1,n+1}\,\setminus *$.





Из изложенного вытекает


\medskip





{\bf Теорема}. {\it Пусть выполнены условия $1$, $2$. Тогда уравнение $(1)$


имеет


решение $x=x(\varepsilon)\to 0$, $\lambda=\lambda(\varepsilon)\to 0$ при


$\varepsilon \to 0$, где


$$


P x(\varepsilon) = \sum^n_{i=1}\varepsilon ^{\alpha_i}(\eta_i^0 +


o(1))\varphi_i,


$$


$\lambda(\varepsilon)\sim \varepsilon^{\alpha_{n+1}}\eta^0_{n+1}. $


Последовательность $\{x_m,\lambda_m\}$, определяемая формулами $(17)$--$(21)$,


сходится к этому решению в окрестности нуля.}





\begin{note}


Правые части формулы (21) содержат отрицательные степени


$\varepsilon$. Но это устранимые особенности. В случае полиномиальных


нелинейностей от них легко избавиться, сократив правые части в (21) на


соответствующие степени $\varepsilon$ (см.\ ниже пример). После этого


сходимость будет равномерной по $\varepsilon$ в окрестности


$|\varepsilon|\le\rho$. Если сокращение нельзя сделать явно, то для


обеспечения устойчивости согласно [17] можно в отрицательных степенях


$\varepsilon$ в правых частях (21) сделать замену $\varepsilon \Rightarrow


\varepsilon +\, \text{sign}\, \varepsilon\sigma^{\nu}$,


где $0<\nu<\frac{1}{2p}$, 


$p = \max\limits_{1\le j\le n}(\Theta_j - \min(\alpha_1,\dotsc,\alpha_n))$,


$\sigma$ --- максимальная абсолютная погрешность вычислений. Тогда


рассмотренный итерационный метод будет регуляризующим алгоритмом в


смысле А.Н.\,Тихонова.


\end{note}





\section{Одноступенчатый вариант схемы (17)--(21)}





{\bf Условие 3.} {\it Пусть диаграммы Ньютона коэффициентов проекции $QF$


имеют параллельные грани $pi+qk = \Theta_j$, $j=1,\dotsc,n;$ $p, q,


\Theta_j \in N^+$,


$R_{0i} = 0$ при $i<\frac {p}{q}$, $QR_{\frac pq}=0$.}





Тогда на основании леммы 3 из [13] условие 1 будет выполнено при $\alpha_1


= \dotsb = \alpha_n = p$, $\alpha_{n+1} = q$.





Положим $\xi_i = \varepsilon^p\eta_i$, $i=1,\dotsc,n$,


$\lambda = \varepsilon^q\eta_{n+1}$. При этом система (8) имеет вид


\begin{equation}


l_j \triangleq


\sum_{pi+qk=\Theta_j}\bigg\la R_{ik}\bigg(\sum^n_{s=1}\eta_s\varphi_s +


x^0(\eta_{n+1})\bigg), \psi_j\bigg\ra\eta^k_{n+1} = 0,\quad


j=\overline{1,n},


\end{equation}


где


$$


x^0 =\begin{cases}


0, & \frac pq \notin N^+; \\


\Gamma R_{0\frac pq}\eta^{\frac pq}_{n+1}, &


\frac pq \in N^+.


\end{cases}


$$


Пусть система (22) имеет решение $\eta^0$ полного ранга.


Тогда условия теоремы будут выполнены. В итерационных формулах (17)--(21)


в этом случае надо положить $N=1, r =\begin{cases}


p, & \text{если } p\le q; \\


q, & \text{eсли } p>q.\end{cases}$


\medskip





{\bf Пример.} Рассмотрим уравнение


\begin{equation}


Bx = \lambda Ax + F_l(x),


\end{equation}


где $A\in L(E_1 \to E_2)$, $l\ge 2$.


Пусть $QA\ne 0$, $QF_l(\xi\varphi\not\equiv 0.)$ Тогда условие 3


выполняется


при $p=1$, $q=l-1$, $\Theta_1 =\dotsb =\Theta_n =l$. Алгебраическая система


(22) примет вид


\begin{equation}


l_j(\eta_1,\dotsc,\eta_{n+1}) \triangleq


\eta_{n+1}\sum^n_{k=1}\la A\varphi_k,\psi_j\ra\eta_k +


\la F_l(\eta\varphi),\psi_j\ra =


 0, \quad j=\overline{1,n}.


\end{equation}





Пусть $\eta^0$ --- решение полного ранга системы (24), причем


$$


D = \bigg\|\frac{\partial l_j(\eta^0)}{\partial \eta_i}


\bigg\|_{i,j=\overline{1,n}},\quad


\det D \ne 0.


$$


Тогда можно использовать схему (17)--(21) при $r=1$, $N=1$, $y^0=0$, $*=n+1$,


$\lambda = \varepsilon^{l-1}\eta^0_{n+1}$.





Соответствующая последовательность $x_m = \varepsilon(\eta^m\varphi +


\wh{x}_m)$, где $\wh{x}_m$ --- единственное решение линейного уравнения


$$


\ov{B}\wh{x}_m = \varepsilon^{l-1}\{\eta ^0_{n+1}A(\eta^{m-1}\varphi +


\wh{x}_{m-1}) + F_l(\eta^{m-1}\varphi + \wh{x}_{m-1})\}


$$


c непрерывно обратимым оператором


$\ov{B} = B + \sum\limits_1^n\la\cdot,\gamma_i\ra z_i$,


и $\eta^m = \eta^{m-1} - \la\eta^0_{n+1}A(\eta^{m-1}\varphi + \wh{x}_m) +


F_l(\eta^{m-1}\varphi + \wh{x}_m), D^{-1}\Psi\ra$, $m=1,2,\dotsc$,


$\eta^0 = (\eta^0_1,\dotsc,\eta^0_n)'$, $\wh{x_0}=0$,


cходится к решению $x(\varepsilon)$ уравнения (23) при


$|\varepsilon|<\rho$.





Пусть в этом примере $E_1=E_2=H$, операторы $B$, $A$ самосопряженные


$\la F_l(\eta\varphi),\varphi_i\ra =


\frac{\partial V}{\partial \eta_i}$, $i=1,\dotsc,n$,


и матрица $\|\la A\varphi_i,\varphi_k\ra\|^n_{i,k=1}$


определенно положительная


(или определенно отрицательная). Тогда гарантировано существование


ненулевых вещественных решений  системы (24). А именно, за


$(\eta^0_1,\dotsc,\eta^0_n)$


можно брать точки экстремумов потенциала $V(\eta)$ на компакте


$\frac12\sum\limits^n_{i,j=1}\la A\varphi_i,\varphi_j\ra\eta_i\eta_j =1$


(или на $\frac 12\sum\limits^n_{i,j=1}


\la A\varphi_i,\varphi_j\ra\eta_i\eta_j = -1)$,


а за $\eta^0_{n+1}$ --- соответствующие множители Лагранжа.





Примеры решения интегральных уравнений одноступенчатым вариантом подобной


схемы см. в [13]--[15].





\begin{note}


В общем случае система (24) может не иметь решений полного ранга.


Соответствующие разветвляющие решения ``неполного ранга'' требуют


привлечения


других методов и пока изучены слабо. В условиях групповой симметрии [18]


уравнения (1) в этом случае для построения итерационных схем можно


привлечь


результаты из [11], [13], [15]. Смотри также п.\,3, пример в [13],


работы [19], [15].


\end{note}





\section*{Заключение}





В процессе итераций (17)--(21) в общем случае строятся три


последовательности


\begin{equation}


\{\varepsilon^r\Gamma y_m\}, \ \; \{\xi^m\varphi\}\ \; \text{и}\ \;


\{\lambda_m\}.


\end{equation}


Первая из них сходится к проекции $(J-P)x(\varepsilon)$, вторая --- к


проекции


$Px(\varepsilon)$ искомого решения, $\lim\limits_{m\to\infty}\lambda_m =


\lambda(\varepsilon)$. В случае $N \ge 2$ при построении $y_m$ требуется


решить $N-1$  линейное уравнение с непрерывно обратимым оператором.


Вектор $\xi^m$ ($\lambda_m$, если $*\ne n+1$) определяется по вычисленному


$\Gamma y_m$.


Таким образом, всего для определения $m$-х членов последовательностей (25)


надо решить $N$ линейных уравнений.





Схема (17)--(21) реализует в итерационном виде развитие идей аналитического


метода Ляпунова--Шмидта [1]. В дополнение к этим идеям допускается смена


параметра униформизации ветвей. Например, наряду с параметром  $\lambda$,


входящим в уравнение, можно использовать любой из коэффициентов проекции


$Px$ искомого решения. Это дает возможность ``обходить'' некоторые типы


точек ветвления и в сочетании с результатами работы [20] расширить область


сходимости метода.








\begin{thebibliography}{99}





\bibitem{1}


Вайнбеpг М.М., Тpеногин В.А. {\em Теоpия ветвления pешений нелинейных


уpавнений}. -- М.: Наука,  1969. -- 527 с.


\bibitem{2}


Брюно А.Д. {\em Степенная геометрия в алгебраических и


дифференциальных уравнениях}. -- М.: Наука, 1998. -- 288 с.


\bibitem{3}


Zeidler E. {\em Nonlinear functional analyses and its Applications}. I:


{\em Fixed-point theorems}. -- Transl. from Germ. New~York:


Springer-Verlag, 1986. -- 897 p.


\bibitem{4}


Тpеногин В.А. {\em Функциональный анализ}. -- М.: Наука, 1980. -- 491 с.


\bibitem{5}


Логинов Б.В. {\em Теория ветвления решений нелинейных уравнений в


условиях групповой симметрии}. -- Ташкент: ФАН, 1985. -- 184 c.


\bibitem{6}


Keller N.B. {\em Numerical solution of bifurcation and nonlinear


problems} // Application of bifurcation theory, % (ed. P.Rabinovich),


New~York: Academic Press, 1977. -- P.\,359--384.


\bibitem{7}


Moore G. {\em The numerical treatment of nontrivial bifurcation


points} //


Numer. funct. anal. and optim. -- 1980. -- \No\,2. -- P.\,44--72.


\bibitem{8}


Allgower E.L., Schwetlick H. {\em A general view of minimally extended


system for simple bifurcation points} // Z. angew. math. mech. -- 1997. --


Bd.\,77. -- \No\,2. -- S.\,83--97.


\bibitem{9}


Langford W.F. {\em Numerical solution of bifurcation problems for


ordinary differential


equation} // Numer. math. -- 1977. -- Bd.\,28. -- \No\,2. -- S.\,71--190.


\bibitem{10}


Сидоpов Н.А. {\em Общие вопpосы pегуляpизации в задачах теоpии


ветвления}. -- Иркутск: Изд-во Иpкутск. ун-та, 1982. -- 312 с.


\bibitem{11}


Логинов Б.В., Сидоpов Н.А. {\em Гpупповая симметpия уpавнения pазветвления


Ляпунова--Шмидта и итеpационные методы в задаче о точке


бифуpкации} // Матем. сб. -- 1991. -- T.\,182. -- \No\,5. -- C.\,681--691.


\bibitem{12}


Сидоpов Н.А. {\em Явная и неявная паpаметpизация пpи постpоении


pазветвляющихся


pешений итеpационными методами} // Матем. сб. -- 1995. -- T.\,186.


-- \No\,2. -- C.\,129--141.


\bibitem{13}


Сидоров Н.А. {\em $N$-ступенчатый итерационный метод в теории


ветвления решений нелинейных уравнений} // Сиб. матем. журн. -- 1997. --


T.\,38. -- \No\,2. -- C.\,383--395.


\bibitem{14}


Ермилова Н.В., Марканова Д.Ю. {\em Итерационные методы


разветвляющихся решений в случае квазилинейной главной части уравнения


разветвления} // В кн.: Приближенные методы анализа. -- Иркутск:


Изд-вo Иркутск. гос. пед. ин-та, 1997. -- C.\,54--63.


\bibitem{15}


Сидоров Н.А., Ермилова Н.В., Марканова Д.Ю. {\em $N$-ступенчатые


итерации, коммутируемость и квазилинейность уравнения разветвления в


методе Ляпунова--Шмидта} // В кн.: Приближенные методы анализа. -- Иркутск:


Изд-вo Иркутск. гос. пед. ин-та, 1997. -- C.\,75--87.


\bibitem{16}


Сидоров Н.А. {\em О неявной параметризации решений системы


разветвления


Ляпунова--Шмидта}. -- Ред. журн. ``Изв. вузов. Математика''.


-- 1998. -- 12 c. -- Деп. ВИНИТИ \No\,2475-B98.


\bibitem{17}


Сидоров Н.А., Треногин В.А. {\em Регуляризация простых решений


нелинейных уравнений в окрестности точки ветвления} // Сиб. матем. журн.


-- 1978. -- T.\,20. -- \No\,1. -- C.\,180--183.


\bibitem{18}


Овсянников Л.В. {\em Гpупповой анализ диффеpенциальных уpавнений}.


-- М.: Наука, 1978. -- 399~с.


\bibitem{19}


Сидоров Н.А. {\em $(T-M)$-коммутируемость, ранг и итерации в теории


разветвляющихся решений} // В кн.: Тр. 11-й Байкальск. международн.


школы-семинара, 5--12 июля, 1998, секция 4. Численный анализ, обратные и


некорректные задачи. Иркутск, 1998. -- C.\,169--172.


\bibitem{20}


Sidorov N.A.,Abdullin V.R. {\em Interlaced branching equations and


invariance in the theory of nonlinear equations} //


In: ``Proc. of the Conf. Symmetry and Perturbation Theory",


Roma, December 16-22, 1998, World Scientific, 1999. -- P.\,309--313.





\end{thebibliography}





\vskip 2cm





\post{\it Иркутский государственный}{\it Поступила}


\post{\it университет}{11.02.2000}





\label{end}


\end{document}


